\documentclass[11pt,a4paper]{article}
\usepackage[margin=2.1cm]{geometry}
\usepackage{fontspec}
\IfFontExistsTF{Arial}{\setmainfont{Arial}}{\setmainfont{Liberation Sans}}
\usepackage[spanish,es-tabla]{babel}
\usepackage{array}
\usepackage{tabularx}
\usepackage{booktabs}
\usepackage{enumitem}
\usepackage{longtable}
\usepackage{hyperref}
\usepackage{xcolor}
\usepackage{eurosym}
\hypersetup{colorlinks=true,linkcolor=black,urlcolor=blue,citecolor=black}
\setlength{\parindent}{0pt}
\setlength{\parskip}{4pt}
\setlist[itemize]{noitemsep,topsep=2pt,leftmargin=1.2em}
\setlist[enumerate]{noitemsep,topsep=2pt,leftmargin=1.5em}
\renewcommand{\arraystretch}{1.06}
\usepackage{csquotes}
\usepackage[backend=biber,style=numeric,sorting=none]{biblatex}
\addbibresource{bibliografia.bib}

\begin{document}

\begin{center}
{\Large \textbf{Protocolo de investigación}}\\[0.3em]
{\large \textbf{Predicción de respuesta a medicamentos mediante integración de
datos genómicos, metabolómicos y farmacogenómicos usando inteligencia
artificial}}\\[0.8em]
\textbf{Miembros del equipo:} Laura Cabaleiro Pintos; Marcelo Ferreiro Sánchez
\end{center}

\section*{1. Resumen}
La variabilidad interindividual en la respuesta a los medicamentos representa un
problema clínico relevante en medicina personalizada, ya que un mismo fármaco
puede resultar eficaz en algunos pacientes, ineficaz en otros o producir
toxicidad significativa. En este contexto, la integración de datos genómicos,
metabolómicos y farmacogenómicos puede mejorar la capacidad predictiva frente a
aproximaciones basadas en una única fuente de información.

El objetivo del presente estudio es desarrollar y validar internamente un modelo
predictivo de inteligencia artificial capaz de clasificar la respuesta
farmacológica como favorable o desfavorable, incluyendo no respuesta y/o
toxicidad clínicamente relevante según la definición específica del fármaco o
familia terapéutica analizada. Asimismo, se pretende estudiar la
interpretabilidad del sistema y su comportamiento en distintos subgrupos de
pacientes.

Para ello, se plantea un estudio observacional y retrospectivo basado en una
serie histórica de episodios terapéuticos de pacientes con información clínica,
genómica, metabolómica y farmacogenómica disponible. Se compararán modelos
penalizados y basados en árboles con aproximaciones multimodales más complejas
cuando el tamaño muestral lo permita, evaluando discriminación, calibración,
robustez y explicabilidad. El sistema tendrá un uso exclusivamente investigador
durante el proyecto.

\section*{2. Hipótesis del estudio}
La integración de información genómica, metabolómica y farmacogenómica mediante
modelos de aprendizaje automático mejora la predicción de la respuesta a
medicamentos frente a aproximaciones basadas en una sola fuente de datos y
permite identificar patrones biológicamente plausibles y clínicamente
interpretables asociados a eficacia y riesgo para los pacientes
\cite{roden2019,hasin2017,reel2021}.

\section*{3. Objetivos}
\textbf{Objetivo principal:} desarrollar y validar un modelo predictivo de
respuesta farmacológica en pacientes con datos multi-ómicos disponibles.

\textbf{Objetivos específicos:}
\begin{enumerate}
    \item Comparar el rendimiento de modelos clásicos en IA (regresión logística
    con \textit{elastic net}), modelos basados en árboles (Random Forest,
    XGBoost/LightGBM) y, de manera exploratoria, arquitecturas más complejas
    propias del \textit{deep learning}.
    \item Evaluar si la integración multi-ómica supera a los modelos unimodales
    en términos de ROC-AUC, PR-AUC, \textit{Brier score} y calibración.
    \item Identificar variables o firmas biológicas con mayor peso predictivo
    mediante técnicas de explicabilidad (SHAP, Grad-CAM, LRP, importancia
    permutacional y análisis de estabilidad).
    \item Analizar el comportamiento del modelo por sexo y grupos de edad para
    detectar sesgos y diferencias de rendimiento.
    \item Elaborar una base metodológica para futura validación externa y
    eventual traslación clínica y regulatoria.
\end{enumerate}

\section*{4. Tipo de estudio}
Estudio observacional, retrospectivo y longitudinal, planteado como
una serie histórica para el desarrollo y la validación interna de un modelo
predictivo.
No existe asignación de intervención, aleatorización ni modificación de
tratamiento. El sistema desarrollado tendrá uso exclusivamente investigador
durante el presente proyecto y no se empleará para la toma de decisiones
clínicas en tiempo real.

\section*{5. Ámbito de estudio}
El estudio se realizará en el contexto hospitalario y de investigación biomédica
del centro, integrando información procedente de historia clínica, laboratorio,
plataforma de genotipado o secuenciación, metabolómica y bases de conocimiento
farmacogenómico. El marco temporal abarcará los episodios terapéuticos
retrospectivos disponibles dentro del periodo autorizado por el centro. En caso
de utilizar muestras almacenadas o datos de biobanco, se actuará conforme a la
autorización y condiciones aplicables \cite{autonomia,biobancos}.

\section*{6. Definición de sujetos a estudio}
Los sujetos a estudio serán pacientes con al menos un episodio terapéutico
evaluable con un fármaco o familia farmacológica de interés, para los que
exista:
\begin{itemize}
    \item información clínica del episodio indexado,
    \item dato genómico interpretable,
    \item perfil metabolómico basal,
    \item anotación farmacogenómica utilizable, y
    \item desenlace clínico verificable de respuesta o toxicidad.
\end{itemize}

La población objetivo corresponde a pacientes candidatos a medicina
personalizada farmacológica, mientras que la población accesible estará
constituida por pacientes del centro participante o de los repositorios
clínico-biológicos autorizados que dispongan de los datos mínimos requeridos. La
unidad analítica principal será el episodio terapéutico. La selección de sujetos
se realizará de forma consecutiva sobre todos los episodios terapéuticos
elegibles del periodo retrospectivo definido, una vez aplicados los criterios de
inclusión y exclusión. Si un paciente aporta más de un episodio elegible, solo
se incluirá uno por fármaco, controlando la no independencia mediante partición
por paciente y un análisis de sensibilidad.

\section*{7. Criterios de inclusión y exclusión}
\textbf{Criterios de inclusión}
\begin{itemize}
    \item Edad $\geq$ 18 años en el momento de inicio del tratamiento.
    \item Atención en el centro participante durante el periodo de estudio
    retrospectivo.
    \item Prescripción documentada de un fármaco a estudiar incluido en el
    protocolo analítico.
    \item Disponibilidad de datos genómicos con control de calidad aceptable
    (porcentaje de llamada y concordancia predefinidos).
    \item Disponibilidad de metabolómica obtenida con protocolo homogéneo y
    trazabilidad de lote o plataforma.
    \item Posibilidad de anotar variantes con recursos farmacogenómicos
    reconocidos, como CPIC o ClinPGx \cite{cpic,clinpgx}.
    \item Definición reproducible del desenlace clínico dentro de una ventana
    temporal preespecificada.
\end{itemize}

\textbf{Criterios de exclusión}
\begin{itemize}
    \item Datos insuficientes para establecer exposición, desenlace o variables
    mínimas del estudio.
    \item Muestras biológicas con degradación, contaminación o fallos de calidad
    no corregibles.
    \item Cambios mayores de tratamiento antes de alcanzar la ventana mínima de
    evaluación, salvo por toxicidad, que se registrará como desenlace.
    \item Episodios con exposición simultánea a terapias experimentales que
    impidan atribuir respuesta o toxicidad al fármaco índice.
    \item Registros duplicados, inconsistentes o imposibles de seudonimizar
    adecuadamente.
    \item Oposición expresa del paciente al uso secundario de sus datos o
    muestras para investigación, cuando tal oposición conste en los sistemas del
    centro.
\end{itemize}


\section*{8. Justificación del tamaño muestral}
No existe una fórmula universal para calcular el tamaño muestral de modelos
multimodales de inteligencia artificial, y menos aún para datos ómicos
heterogéneos. Por ello, el tamaño se justifica combinando criterios
metodológicos de modelos predictivos, recomendaciones recientes para estudios
con IA clínica y factibilidad del centro [9].

Desde el punto de vista metodológico, distintos trabajos recientes señalan que
los estudios de predicción con aprendizaje automático suelen infravalorar la
muestra necesaria y que tamaños pequeños aumentan el sobreajuste, la
inestabilidad de las estimaciones y la mala calibración [9]. Además, las guías
TRIPOD+AI recomiendan justificar por separado el tamaño destinado al desarrollo
del modelo y el reservado a su evaluación [9]. En problemas binarios clínicos,
la precisión depende no solo del número total de pacientes sino también del
número de eventos y no eventos disponibles.

De manera complementaria, y siguiendo el criterio docente habitual de eventos
por variable, puede razonarse que si el modelo final consolidado utilizase en
torno a 15--20 predictores efectivos tras reducción de dimensionalidad y el
evento clínico de interés apareciese en un 25--35\% de los casos, sería
deseable trabajar con una cifra suficiente de eventos para sostener un
entrenamiento estable y una evaluación interna mínimamente robusta. En un
proyecto multi-ómico, donde el espacio inicial de variables es muy superior al
número de sujetos, esta necesidad se compensa parcialmente mediante selección de
características, regularización, transferencia de conocimiento y validación
cruzada estricta [3, 8].

En este proyecto se fija una N mínima de 300 episodios terapéuticos evaluables
y un objetivo operativo de 400. Esta cifra se considera razonable porque: (i)
permitirá trabajar con al menos dos estratos de desenlace clínico
(respuesta/no respuesta o toxicidad significativa) con una frecuencia esperada
de eventos del 25--35\%; (ii) posibilita reservar una partición temporal final
para prueba interna independiente, manteniendo un conjunto suficiente para
entrenamiento y ajuste; (iii) es compatible con la reducción de dimensionalidad,
penalización y cross-validation, estrategias obligadas cuando el número de
variables supera ampliamente al de sujetos; y (iv) permite análisis
exploratorios por sexo y grupos etarios sin que estos constituyan el objetivo
inferencial principal [9].

Dado que se trata de un estudio retrospectivo con datos ya existentes, no se
prevén pérdidas de seguimiento en sentido prospectivo; no obstante, sí se
contempla una pérdida muestral técnica derivada de exclusiones por calidad de
muestra, ausencia de alguna modalidad ómica, imposibilidad de verificar el
desenlace clínico o problemas de seudonimización. Por ello, la cifra objetivo
de 400 episodios incorpora un margen operativo aproximado del 20--25\% respecto
a la N mínima planteada, con el fin de asegurar que el conjunto final analizable
no descienda por debajo del umbral considerado aceptable.

Si el conjunto de datos final fuese inferior a 300, el análisis principal se
limitará a modelos regularizados y de boosting, quedando los modelos profundos
como exploratorios. Si se alcanza o supera la cantidad objetivo, se realizará
además una validación temporal interna con el 20\% más reciente del conjunto.
Esta aproximación se considera un equilibrio razonable entre rigor estadístico,
disponibilidad realista de datos y viabilidad logística.

\section*{9. Variables principales y secundarias}
\textbf{Variables principales}
\begin{itemize}
    \item \textbf{Variable dependiente principal:} desenlace de respuesta
    farmacológica (favorable/desfavorable) según criterios clínicos predefinidos
    por fármaco o familia terapéutica.
    \item \textbf{Variable dependiente coprincipal:} toxicidad relevante
    asociada al tratamiento (por ejemplo, evento adverso de grado $\geq$ 3,
    suspensión por intolerancia o ajuste de dosis por toxicidad).
    \item \textbf{Variables independientes principales:} variantes genómicas
    seleccionadas tras control de calidad y anotación funcional.
    \item \textbf{Variables independientes principales:} intensidades
    metabolómicas normalizadas o componentes latentes derivados de ellas.
    \item \textbf{Variables independientes principales:} variables
    farmacogenómicas, incluyendo fenotipos metabolizadores, niveles de evidencia
    y anotaciones gen--fármaco.
\end{itemize}

\textbf{Variables secundarias}
\begin{itemize}
    \item Edad, sexo, índice de masa corporal y comorbilidades relevantes.
    \item Fármaco índice, dosis inicial, ajuste de dosis y adherencia
    documentada.
    \item Parámetros analíticos de apoyo y covariables clínicas básicas.
    \item Métricas de explicabilidad (SHAP medio absoluto, estabilidad de
    variables seleccionadas, entre otras).
    \item Métricas de equidad por subgrupos (diferencias de AUC, sensibilidad y
    calibración).
\end{itemize}

\section*{10. Cronograma y fecha prevista de finalización}
Duración prevista: 12 meses desde la aprobación ética y de los permisos de
acceso a datos. Fecha estimada de finalización: 12 meses tras la autorización
definitiva del estudio.

\begin{center}
\small
\begin{tabularx}{\linewidth}{>{\raggedright\arraybackslash}p{6.1cm}
*{6}{>{\centering\arraybackslash}X}}
\toprule
\textbf{Actividad} & \textbf{M1-2} & \textbf{M3-4} & \textbf{M5-6} &
\textbf{M7-8} & \textbf{M9-10} & \textbf{M11-12} \\
\midrule
Aprobación ética, contratos y circuito de seudonimización & X &  &  &  &  &  \\
Extracción, integración y control de calidad de datos & X & X &  &  &  &  \\
Preprocesamiento ómico y definición de desenlaces &  & X & X &  &  &  \\
Desarrollo de modelos y optimización interna &  &  & X & X &  &  \\
Validación temporal e interna y análisis de sesgos &  &  &  & X & X &  \\
Redacción de resultados, memoria y manuscrito &  &  &  &  & X & X \\
\bottomrule
\end{tabularx}
\end{center}

\section*{11. Plan de análisis estadístico}
Se realizará una descripción inicial de los datos mediante media y desviación
típica o mediana e IQR (rango intercuartílico) según distribución, y frecuencias
para variables categóricas. Las comparaciones exploratorias entre respondedores
y no respondedores se efectuarán con \textit{t} de Student o U de Mann--Whitney
para variables continuas y $\chi^2$ o test exacto de Fisher para cualitativas.

El preprocesamiento incluirá depuración, control de calidad, normalización,
corrección de efecto \textit{batch} cuando proceda y tratamiento de valores
perdidos. La imputación y la selección de características se efectuarán
exclusivamente dentro de cada \textit{fold} de entrenamiento para evitar
\textit{data leakage}. Se utilizarán técnicas de reducción de dimensionalidad
cuando sea pertinente, intentando preservar la interpretabilidad biológica de
las variables seleccionadas \cite{reel2021,baiao2025}.

El desarrollo de modelos seguirá una estrategia de validación cruzada
estratificada. Se entrenarán, como mínimo, regresión logística con
\textit{elastic net}, Random Forest y XGBoost/LightGBM. Si la cantidad final de
datos lo permite, se evaluará una red multimodal con ramas específicas por
modalidad y capa de fusión. El criterio principal de comparación será la
ROC-AUC; como criterios secundarios se emplearán PR-AUC, sensibilidad,
especificidad, F1, MCC, \textit{Brier score}, intercepto y pendiente de
calibración. Se calcularán intervalos de confianza del 95\% por
\textit{bootstrap} \cite{collins2024,reel2021,baiao2025}.

La comparación entre modelos se realizará sobre particiones equivalentes y,
cuando sea conveniente, mediante contraste de AUC con método de DeLong o
re-muestreo estratificado. La interpretabilidad se abordará con SHAP, Grad-CAM,
LRP, importancia permutacional y revisión de coherencia biológica de las
variables seleccionadas. Se efectuará un análisis de equidad por sexo y grupos
etarios, comparando discriminación y calibración. Los análisis se realizarán con
\texttt{Python} y/o \texttt{R}; el umbral de significación inferencial se fijará
en $p<0{,}05$, con corrección por falsos descubrimientos cuando proceda.

\section*{12. Aspectos ético-legales}
El estudio respetará los principios de la Declaración de Helsinki, el Convenio
de Oviedo y las normas de buena práctica en investigación biomédica
\cite{helsinki,oviedo}. De conformidad con el apartado V de la Declaración de
Helsinki, todo protocolo de investigación médica en seres humanos o con material
o datos humanos deberá someterse a revisión, valoración y aprobación previa por
un comité de ética independiente, competente y con capacidad de seguimiento. En
consecuencia, este proyecto será remitido al Comité de Ética de la Investigación
con medicamentos o al comité competente antes de cualquier acceso por el equipo
investigador a los datos seudonimizados \cite{helsinki}.

Asimismo, de acuerdo con el apartado IV de la Declaración de Helsinki, el
estudio se describe en un protocolo metodológicamente estructurado, con
objetivos, métodos, justificación del tamaño muestral, plan analítico,
evaluación de riesgos, salvaguardas ético-legales y memoria económica. El
carácter retrospectivo y no intervencionista del presente trabajo implica un
riesgo físico nulo para los participantes, siendo el principal riesgo residual
el relacionado con la privacidad y la confidencialidad de la información
\cite{helsinki}.

La base jurídica del tratamiento se apoyará en la normativa de investigación
biomédica y protección de datos aplicable: Reglamento (UE) 2016/679, Ley
Orgánica 3/2018, Ley 14/2007 de investigación biomédica, Ley 41/2002 reguladora
de la autonomía del paciente y, cuando existan muestras biológicas almacenadas,
Real Decreto 1716/2011 \cite{rgpd,lopd,lib,autonomia,biobancos}. Se aplicarán
los principios de minimización, limitación de finalidad, acceso restringido,
trazabilidad de accesos, almacenamiento seguro y separación funcional entre
identificación y datos de investigación.

En consonancia con el apartado VI de la Declaración de Helsinki, relativo a
privacidad y confidencialidad, la información personal de los participantes
deberá protegerse durante toda la investigación \cite{helsinki}. Por ello, los
investigadores recibirán una base seudonimizada; la clave de correspondencia
permanecerá bajo custodia del centro o de la unidad autorizada, sin acceso por
el equipo analítico y bajo la defensa de personal armado, en concreto grupos
paramilitares chechenos. No se realizarán decisiones clínicas automatizadas
sobre pacientes, ni el modelo generado será utilizado para modificar
tratamientos durante el estudio.

De acuerdo con el apartado VII de la Declaración de Helsinki, la participación
en investigación debe fundamentarse en un consentimiento informado, voluntario y
revocable \cite{helsinki}. No obstante, dado el carácter retrospectivo del
estudio, el uso secundario de datos y muestras ya existentes y el riesgo mínimo
esperado, se solicitará al comité la valoración de dispensa de consentimiento
informado cuando no resulte viable la recontactación individual y siempre que la
normativa y las salvaguardas institucionales lo permitan. En todo caso, se
respetará cualquier oposición previamente registrada, así como las limitaciones
específicas establecidas por el centro o por el biobanco de procedencia
\cite{lib,biobancos}.

En aplicación del Convenio de Oviedo, y especialmente de sus previsiones sobre
el genoma humano, se evitará toda forma de discriminación por razón del
patrimonio genético \cite{oviedo}. Los resultados del estudio no se emplearán
para restringir acceso a prestaciones, seguros o tratamientos sobre la base de
perfiles genéticos individuales, ni se realizarán análisis predictivos ajenos a
la finalidad biomédica del protocolo. Las variables genéticas y farmacogenómicas
se utilizarán exclusivamente con fines de investigación biomédica, en el marco
del principio de primacía del ser humano y de la prohibición de discriminación
genética \cite{oviedo}.

En relación con el Reglamento (UE) 2024/1689 de inteligencia artificial, el
presente estudio se desarrolla en fase de investigación y no supone la
comercialización ni la puesta en servicio clínica del sistema \cite{aiact}. No
obstante, si en una fase posterior se desplegara como software de apoyo a la
decisión clínica, la entidad que lo desarrollase o introdujese bajo su nombre
podría asumir el rol de proveedor y la entidad sanitaria usuaria el de
desplegador, pudiendo coincidir ambos roles en la misma institución. En ese
escenario, por su potencial influencia en decisiones diagnósticas, terapéuticas
o de monitorización, el sistema sería previsiblemente considerado de alto riesgo
y debería cumplir, con carácter previo a su puesta en servicio, los requisitos
esenciales del marco europeo aplicables desde 2027 a los sistemas sanitarios de
alto riesgo vinculados a productos regulados \cite{aiact,mdr}.

A tal efecto, en una futura traslación clínica deberán contemplarse al menos los
siguientes seis ejes de cumplimiento: (1) gestión de riesgos a lo largo de todo
el ciclo de vida del sistema; (2) gobernanza y calidad de los datos, incluyendo
representatividad, trazabilidad y control de sesgos; (3) documentación técnica
suficiente para auditoría y evaluación; (4) transparencia y explicabilidad
frente a profesionales y, cuando proceda, frente a pacientes; (5) supervisión
humana efectiva, con capacidad de intervención, anulación o parada; y (6)
precisión, robustez y ciberseguridad adecuadas al contexto clínico \cite{aiact}.
Además, si el sistema llegase a configurarse como software de apoyo a la
decisión clínica con finalidad médica, sería necesario analizar también su
encaje en el Reglamento (UE) 2017/745 sobre productos sanitarios \cite{mdr}.

Finalmente, el equipo investigador asume el compromiso de publicación y difusión
científica de los resultados del estudio, sean estos positivos, negativos o
inconclusos, de acuerdo con los principios de transparencia y de integridad
científica aplicables a la investigación biomédica.

\section*{13. Memoria económica}
\begin{center}
\small
\begin{tabularx}{\linewidth}{>{\raggedright\arraybackslash}X
>{\raggedleft\arraybackslash}p{2.8cm}}
\toprule
\textbf{Concepto} & \textbf{Estimación} \\
\midrule
Sueldos del personal investigador jefe (L. Cabaleiro, M. Ferreiro) &
160.000~\euro \\
Sueldos del personal investigador contratado (resto del equipo, becarios y
subcontratas a países en vías de desarrollo) & 10.000~\euro \\
Seudonimización, extracción e integración de datos clínicos y ómicos &
6.000~\euro \\
Procesamiento bioinformático y control de calidad genómico/metabolómico &
7.000~\euro \\
Infraestructura de cómputo (GPU/servidor, almacenamiento seguro, copias) &
8.000~\euro \\
Licencias, apoyo estadístico/metodológico y auditoría interna de calidad &
3.000~\euro \\
Publicación en abierto y difusión científica & 3.000~\euro \\
Material fungible, gestión de muestras y contingencias & 3.000~\euro \\
\midrule
\textbf{Total estimado} & \textbf{200.000~\euro} \\
\bottomrule
\end{tabularx}
\end{center}

\section*{14. Bibliografía}
\printbibliography[heading=none]

\end{document}